\section{Theory: What the Instrument Measures}
\label{sec:theory}

A benchmark earns trust when its central metric is justified, its evaluator cannot be
gamed, and its sample sizes are principled, so we treat theory as part of the instrument.
All results are stated under explicit assumptions and carry honest proof-status labels;
none claims that any real defense is secure or provably degrades. \emph{Every result we
claim is proved in full in this section.} What the Supplementary Material adds are the
derivations of the \emph{classical} propagation results we restate in our notation
(Granovetter/Watts/Kempe), not proofs of our own claims.

\begin{assumption}[Trusted checking]
\label{as:tcb}
The harness records a faithful global trace $\Trace$, and the checker $\Checker$
evaluates the fixed predicates $\Predsec,\Predutil$ as functions of $\Trace$ \emph{only}
(\Cref{def:predicates}).
\end{assumption}
\begin{assumption}[i.i.d.\ episodes]
\label{as:iid}
For fixed $\config$, the $M$ episode traces are independent draws from
$\mathcal{D}(\config)$.
\end{assumption}
\begin{assumption}[Bounded indicators]
\label{as:bounded}
$\Predutil,\Predsec\in\{0,1\}$.
\end{assumption}
\begin{assumption}[Monotone propagation model]
\label{as:prop}
Compromise spreads over $\Gtop$ via a monotone rule $\mathcal{P}$ (\Cref{def:propagation}):
under \emph{independent cascade}, a compromised agent independently compromises each
susceptible out-neighbour with probability $\pinf$; under \emph{threshold}, a susceptible
agent is compromised once at least $\thresh$ in-neighbours are. This is a modelling choice
\defname{} makes so propagation is measurable, not an empirical law; its realism is itself
measured via $\CPR$.
\end{assumption}
\begin{assumption}[Continuity of the aggregator]
\label{as:cont}
The utility--security aggregator $g:[0,1]^2\to[0,1]$ is continuous.
\end{assumption}

\subsection{Characterization of the system-level metric}
\label{sec:theory-metric}
The product form of $\NRPs$ is forced by the desiderata of \Cref{sec:problem}, phrased as
five axioms (D1--D5): boundedness, consistency, monotonicity, homogeneity (each in both
arguments), and reduction---the operative instance is stated and proven in \Cref{thm:metric}.
Under them, $g$ reduces to ASB's metric at one agent.

\begin{proposition}[Reduction to ASB]
\label{prop:reduction}
If $N=1$ and $\Predsec,\Predutil$ are the single-agent security and benign predicates,
then $\ASRs=\mathrm{ASR}$, $\PNAs=\mathrm{PNA}$, and $\NRPs=\mathrm{NRP}$ of
\Cref{eq:asb-nrp}. \emph{Proof in the Supplementary Material.}
\end{proposition}

\begin{theorem}[Axiomatic characterization of $\NRPs$; complete proof]
\label{thm:metric}
Write $\NRPs=g(u,\sigma)$ with utility $u=\PNAs\in[0,1]$ and security level
$\sigma=1-\ASRs\in[0,1]$. Under \Cref{as:cont}, suppose $g$ satisfies:
\emph{(D2) consistency}, $g(u,1)=u$ and $g(0,\sigma)=0$;
\emph{(D4a) utility-homogeneity}, $g(\lambda u,\sigma)=\lambda\,g(u,\sigma)$ for
$\lambda\in[0,1]$; and
\emph{(D4b) security-homogeneity}, $g(u,\lambda\sigma)=\lambda\,g(u,\sigma)$ for
$\lambda\in[0,1]$.
Then $g(u,\sigma)=u\,\sigma$ is the unique such aggregator, satisfying boundedness (D1)
and monotonicity (D3).
\end{theorem}

\begin{proof}[Proof of \Cref{thm:metric}]
By D4a with base $u=1$ and D4b with base $\sigma=1$,
$g(u,\sigma)=u\,g(1,\sigma)=u\sigma\,g(1,1)=u\sigma$ by D2. The boundary $g(0,\sigma)=0$ holds
since $0\cdot\sigma=0$; and $u\sigma$ is bounded in $[0,1]$ and non-decreasing in each
argument. Continuity (\Cref{as:cont}) extends the homogeneity identities from rationals to
all of $[0,1]$; the identities hold pointwise.
\end{proof}

\begin{remark}
\label{rem:metric}
\Cref{thm:metric} singles out the product \emph{relative to D2/D4}: replacing D4 by a
worst-coordinate axiom yields $\min(u,\sigma)$, and weighted log-linearity yields
$u^{w}\sigma^{1-w}$. \defname{} reports $u\sigma$ by default (it generalizes ASB,
\Cref{prop:reduction}) and exposes alternatives as configurable.
\end{remark}

\subsection{Evaluator integrity}
\label{sec:theory-integrity}
We formalize why deterministic checking cannot be hijacked by a successful injection.

\begin{theorem}[Evaluator integrity; complete proof]
\label{thm:integrity}
Under \Cref{as:tcb}, for any two episodes with equal recorded traces $\Trace=\Trace'$,
$\Checker(\Trace)=\Checker(\Trace')$. Consequently, an adversarial manipulation that
changes message content but leaves the recorded security-relevant state unchanged cannot
change $\Predsec(\Trace)$; hence neither $\widehat{\ASRs}$ nor $\widehat{\NRPs}$ can be
inflated or deflated by message content alone.
\end{theorem}

\begin{proof}[Proof of \Cref{thm:integrity}]
By \Cref{as:tcb} the predicates are functions of $\Trace$ only, so equal traces give equal
verdicts. For any manipulation $\pi$ inducing $\Trace_\pi=\Trace$ the verdict
$\Predsec(\Trace_\pi)=\Predsec(\Trace)$, and hence so are
$\widehat{\ASRs},\widehat{\PNAs},\widehat{\NRPs}$.
\end{proof}

\begin{proposition}[An LLM judge is hijackable; complete proof of existence]
\label{prop:judgehijack}
There exist an LLM judge $\Judge$ that reads messages and a message manipulation $\pi$
with $\Trace_\pi=\Trace$ such that
$\Judge(\Trace_\pi,\mathrm{msgs}_\pi)\neq\Judge(\Trace,\mathrm{msgs})$.
\emph{Proof in the Supplementary Material.}
\end{proposition}

\begin{remark}
\Cref{thm:integrity} shows the deterministic verdict cannot be \emph{gamed} through
messages; it does \emph{not} show it is \emph{correct} (a mis-specified predicate can be
systematically wrong). Viewed as a test of the compromise hypothesis, $\Checker$ is a
deterministic decision rule governed by classical testing
theory~\cite{neyman1933problem,lehmann2005testing}; the LLM judge additionally exposes
that decision to manipulation through its message input, as documented for LLM
judges~\cite{zheng2023judging,wang2023fair}. \defname{} keeps $\Checker$ on the critical
path and measures the judge-hijack gap as an ablation (RQ2).
\end{remark}

\subsection{Compromise propagation across topologies}
\label{sec:theory-prop}
Under the independent-cascade instance of \Cref{as:prop}, expected reach depends sharply on
topology. We \emph{restate} four classical results in our notation, and derive from them
the one claim we need. On a chain \emph{seeded at an endpoint}, $a_{1+j}$ is compromised
iff all $j$ intervening edges fire, so the expected additional reach is
$\sum_{j=1}^{n-1}\pinf^{j}\le\pinf/(1-\pinf)$, bounded independently of $n$. On a star, a
centre seed gives reach $\pinf(N-1)$ and $\CPR=\pinf$, whereas a leaf seed gives
$\pinf+\pinf^{2}(N-2)$; the gap is $\pinf(1-\pinf)(n-2)=\Theta(\pinf(1-\pinf)n)$. On a tree
of branching factor $b$ seeded at the root the expected level-$\ell$ count is
$(b\pinf)^{\ell}$, so reach is bounded in depth for $b\pinf<1$ and $\Theta((b\pinf)^{d})$
for $b\pinf>1$: a phase transition at $b\pinf=1$. On the complete graph $K_n$ with $s$
seeds, one round already compromises $s+(n-s)\bigl(1-(1-\pinf)^{s}\bigr)$ nodes, which
tends to all of $V$ as $s$ grows. Under threshold contagion an agent with in-degree below
$\thresh$ is never compromised, so a chain with $\thresh\ge2$ has $\CPR=0$. These are the
classical spread results~\cite{granovetter1978threshold,watts2002simple,%
kempe2003maximizing,chalupa1979bootstrap}---we reproduce the derivations in the
Supplementary Material for completeness rather than claim them---and our own consequence
follows.

\begin{theorem}[Collusion advantage on the star]
\label{thm:collusion}
Fix the independent-cascade model on a star of $N$ nodes with a budget of one seed. A
colluding adversary that places its seed at the centre achieves expected compromise
$C=1+\pinf(N-1)$; an independent adversary placing its seed uniformly at random achieves
expected compromise $\tfrac1N C+\tfrac{N-1}{N}L$, writing $C$ as above and
$L=1+\pinf+\pinf^{2}(N-2)$. Their difference
$\Delta_{\mathrm{coll}}=\pinf(1-\pinf)\tfrac{(N-1)(N-2)}{N}$ is positive,
$\Theta(\pinf(1-\pinf)N)$, for all $\pinf\in(0,1)$, $N\ge3$: coordinated placement
strictly dominates. \emph{Proof in the Supplementary Material}; the quantity proved positive is a
difference in expected compromised \emph{agents}, not the probability difference of
\Cref{def:csr}.}
\end{theorem}

\Cref{thm:collusion} describes the \emph{environment's} propagation model, not any real
defense; its value is that ``collusion on/off'' becomes well posed and predicts where to
look---hub-centric topologies---for degradation, which \Cref{sec:experiments} then tests.

\subsection{Forced coordination}
\label{sec:theory-forced}
\begin{proposition}[Forced-coordination lower bound]
\label{prop:forced}
Let a task be $k$-forced-coordination (\Cref{def:forced}) with witness family
$\mathcal{W}$. Then: (a) no set $S$ of agents with $|S|<k$ can satisfy $\Predutil$; and
(b) if $\Advset$ is a transversal of $\mathcal{W}$, the adversary can force $\Predutil=0$
by having its member in each witness withhold or sabotage its contribution; conversely,
if some witness is adversary-free, the honest agents can complete the task.
\emph{Proof in the Supplementary Material.}
\end{proposition}

\subsection{Sample complexity of the estimators}
\label{sec:theory-sample}
The estimates require bounded episodes, driving the cost analysis. The tools are elementary
concentration inequalities~\cite{boucheron2013concentration,hoeffding1963probability}:
Hoeffding's inequality applies to the means here, and McDiarmid's
inequality~\cite{mcdiarmid1989bounded} handles the non-mean estimators. What the resulting
budget does and does not resolve in the experiments is worked out in
the Supplementary Material.

\begin{theorem}[Concentration and sample complexity]
\label{thm:sample}
Under \Cref{as:iid,as:bounded}, for any $\epsilon,\alpha\in(0,1)$: (a) for each
estimator $\hat\theta\in\{\widehat{\ASRs},\widehat{\PNAs}\}$ with target $\theta$,
$\Prob[|\hat\theta-\theta|\ge\epsilon]\le 2e^{-2M\epsilon^{2}}$, so $M\ge\frac{1}{2\epsilon^{2}}
\ln\frac{2}{\alpha}$ episodes give $\epsilon$-accuracy at confidence $1-\alpha$;
(b) the product is Lipschitz; hence, by a union bound, $\Prob[|\widehat{\NRPs}-\NRPs|\ge
2\epsilon]\le 4e^{-2M\epsilon^{2}}$;
and (c) across a sweep of $G$ configurations, Bonferroni control~\cite{benjamini1995controlling}
gives $M=O(\epsilon^{-2}\ln(G/\alpha))$ per configuration, for $\Theta(GMNH)$ total agent
invocations. \emph{Proof in the Supplementary Material.}
\end{theorem}

\begin{corollary}[Sweep budget]
\label{cor:budget}
Resolving a target $\NRPs$ gap of size $\Delta$ between any two of $G$ swept
configurations with family-wise confidence $1-\alpha$ costs
$O\bigl(GNH\Delta^{-2}\ln(G/\alpha)\bigr)$ agent invocations (a sufficient budget, not a
tight rate). \emph{Proof in the Supplementary Material.}
\end{corollary}

